% Fragmento temporário para a seção ``Espaços Amostrais Discretos''.
% Pode ser inserido no arquivo principal com: % Fragmento temporário para a seção ``Espaços Amostrais Discretos''.
% Pode ser inserido no arquivo principal com: % Fragmento temporário para a seção ``Espaços Amostrais Discretos''.
% Pode ser inserido no arquivo principal com: % Fragmento temporário para a seção ``Espaços Amostrais Discretos''.
% Pode ser inserido no arquivo principal com: \input{medidas-gibbs}

\begin{example}[Medidas de Gibbs]\label{ex:gibbs-measure}
  Seja \(\Omega\) um conjunto finito e não vazio, e considere o espaço de
  eventos \(\mathcal{F} = \mathscr{P}(\Omega)\). Seja também
  \[
    H : \Omega \longrightarrow \mathbb{R}
  \]
  uma função. Em mecânica estatística, a função \(H\) é chamada de
  \emph{Hamiltoniano} e atribui uma energia a cada estado \(\omega \in
  \Omega\). Fixado um parâmetro \(\beta > 0\), que interpretamos como o
  inverso da temperatura (absorvendo a constante de Boltzmann na definição de
  \(\beta\)), definimos a \emph{função de partição} por
  \[
    Z(\beta) = \sum_{\omega \in \Omega} e^{-\beta H(\omega)}.
  \]

  Como \(\Omega\) é finito e cada termo da soma é positivo,
  \(0 < Z(\beta) < \infty\). A \emph{medida de Gibbs} associada a \(H\) e
  \(\beta\) é a medida de probabilidade \(\mathbb{P}_\beta\) definida por
  \[
    \mathbb{P}_\beta(A)
    = \frac{1}{Z(\beta)}
      \sum_{\omega \in A} e^{-\beta H(\omega)},
    \qquad A \in \mathcal{F}.
  \]
  Assim, estados de menor energia recebem maior peso, enquanto a função de
  partição normaliza esses pesos para que a probabilidade total seja 1.
\end{example}

\begin{exercise}\label{exercise:gibbs-high-temperature}
  No contexto do \Cref{ex:gibbs-measure}, fixe um evento \(A \in \mathcal{F}\).
  Calcule o limite de alta temperatura
  \[
    \lim_{\beta \downarrow 0} \mathbb{P}_\beta(A).
  \]
  Interprete o resultado em termos da medida uniforme sobre \(\Omega\).
\end{exercise}

\begin{exercise}\label{exercise:gibbs-low-temperature}
  No contexto do \Cref{ex:gibbs-measure}, defina
  \[
    h_{\min} = \min_{\omega \in \Omega} H(\omega),
    \qquad
    \Omega_{\min} = \{\omega \in \Omega : H(\omega) = h_{\min}\}.
  \]
  Calcule o limite de baixa temperatura
  \[
    \lim_{\beta \to \infty} \mathbb{P}_\beta(A).
  \]
  Mostre que, no limite, a medida se torna uniforme sobre os estados de menor
  energia. Em particular, descreva o que acontece quando o estado de menor
  energia é único.
\end{exercise}


\begin{example}[Medidas de Gibbs]\label{ex:gibbs-measure}
  Seja \(\Omega\) um conjunto finito e não vazio, e considere o espaço de
  eventos \(\mathcal{F} = \mathscr{P}(\Omega)\). Seja também
  \[
    H : \Omega \longrightarrow \mathbb{R}
  \]
  uma função. Em mecânica estatística, a função \(H\) é chamada de
  \emph{Hamiltoniano} e atribui uma energia a cada estado \(\omega \in
  \Omega\). Fixado um parâmetro \(\beta > 0\), que interpretamos como o
  inverso da temperatura (absorvendo a constante de Boltzmann na definição de
  \(\beta\)), definimos a \emph{função de partição} por
  \[
    Z(\beta) = \sum_{\omega \in \Omega} e^{-\beta H(\omega)}.
  \]

  Como \(\Omega\) é finito e cada termo da soma é positivo,
  \(0 < Z(\beta) < \infty\). A \emph{medida de Gibbs} associada a \(H\) e
  \(\beta\) é a medida de probabilidade \(\mathbb{P}_\beta\) definida por
  \[
    \mathbb{P}_\beta(A)
    = \frac{1}{Z(\beta)}
      \sum_{\omega \in A} e^{-\beta H(\omega)},
    \qquad A \in \mathcal{F}.
  \]
  Assim, estados de menor energia recebem maior peso, enquanto a função de
  partição normaliza esses pesos para que a probabilidade total seja 1.
\end{example}

\begin{exercise}\label{exercise:gibbs-high-temperature}
  No contexto do \Cref{ex:gibbs-measure}, fixe um evento \(A \in \mathcal{F}\).
  Calcule o limite de alta temperatura
  \[
    \lim_{\beta \downarrow 0} \mathbb{P}_\beta(A).
  \]
  Interprete o resultado em termos da medida uniforme sobre \(\Omega\).
\end{exercise}

\begin{exercise}\label{exercise:gibbs-low-temperature}
  No contexto do \Cref{ex:gibbs-measure}, defina
  \[
    h_{\min} = \min_{\omega \in \Omega} H(\omega),
    \qquad
    \Omega_{\min} = \{\omega \in \Omega : H(\omega) = h_{\min}\}.
  \]
  Calcule o limite de baixa temperatura
  \[
    \lim_{\beta \to \infty} \mathbb{P}_\beta(A).
  \]
  Mostre que, no limite, a medida se torna uniforme sobre os estados de menor
  energia. Em particular, descreva o que acontece quando o estado de menor
  energia é único.
\end{exercise}


\begin{example}[Medidas de Gibbs]\label{ex:gibbs-measure}
  Seja \(\Omega\) um conjunto finito e não vazio, e considere o espaço de
  eventos \(\mathcal{F} = \mathscr{P}(\Omega)\). Seja também
  \[
    H : \Omega \longrightarrow \mathbb{R}
  \]
  uma função. Em mecânica estatística, a função \(H\) é chamada de
  \emph{Hamiltoniano} e atribui uma energia a cada estado \(\omega \in
  \Omega\). Fixado um parâmetro \(\beta > 0\), que interpretamos como o
  inverso da temperatura (absorvendo a constante de Boltzmann na definição de
  \(\beta\)), definimos a \emph{função de partição} por
  \[
    Z(\beta) = \sum_{\omega \in \Omega} e^{-\beta H(\omega)}.
  \]

  Como \(\Omega\) é finito e cada termo da soma é positivo,
  \(0 < Z(\beta) < \infty\). A \emph{medida de Gibbs} associada a \(H\) e
  \(\beta\) é a medida de probabilidade \(\mathbb{P}_\beta\) definida por
  \[
    \mathbb{P}_\beta(A)
    = \frac{1}{Z(\beta)}
      \sum_{\omega \in A} e^{-\beta H(\omega)},
    \qquad A \in \mathcal{F}.
  \]
  Assim, estados de menor energia recebem maior peso, enquanto a função de
  partição normaliza esses pesos para que a probabilidade total seja 1.
\end{example}

\begin{exercise}\label{exercise:gibbs-high-temperature}
  No contexto do \Cref{ex:gibbs-measure}, fixe um evento \(A \in \mathcal{F}\).
  Calcule o limite de alta temperatura
  \[
    \lim_{\beta \downarrow 0} \mathbb{P}_\beta(A).
  \]
  Interprete o resultado em termos da medida uniforme sobre \(\Omega\).
\end{exercise}

\begin{exercise}\label{exercise:gibbs-low-temperature}
  No contexto do \Cref{ex:gibbs-measure}, defina
  \[
    h_{\min} = \min_{\omega \in \Omega} H(\omega),
    \qquad
    \Omega_{\min} = \{\omega \in \Omega : H(\omega) = h_{\min}\}.
  \]
  Calcule o limite de baixa temperatura
  \[
    \lim_{\beta \to \infty} \mathbb{P}_\beta(A).
  \]
  Mostre que, no limite, a medida se torna uniforme sobre os estados de menor
  energia. Em particular, descreva o que acontece quando o estado de menor
  energia é único.
\end{exercise}


\begin{example}[Medidas de Gibbs]\label{ex:gibbs-measure}
  Seja \(\Omega\) um conjunto finito e não vazio, e considere o espaço de
  eventos \(\mathcal{F} = \mathscr{P}(\Omega)\). Seja também
  \[
    H : \Omega \longrightarrow \mathbb{R}
  \]
  uma função. Em mecânica estatística, a função \(H\) é chamada de
  \emph{Hamiltoniano} e atribui uma energia a cada estado \(\omega \in
  \Omega\). Fixado um parâmetro \(\beta > 0\), que interpretamos como o
  inverso da temperatura (absorvendo a constante de Boltzmann na definição de
  \(\beta\)), definimos a \emph{função de partição} por
  \[
    Z(\beta) = \sum_{\omega \in \Omega} e^{-\beta H(\omega)}.
  \]

  Como \(\Omega\) é finito e cada termo da soma é positivo,
  \(0 < Z(\beta) < \infty\). A \emph{medida de Gibbs} associada a \(H\) e
  \(\beta\) é a medida de probabilidade \(\mathbb{P}_\beta\) definida por
  \[
    \mathbb{P}_\beta(A)
    = \frac{1}{Z(\beta)}
      \sum_{\omega \in A} e^{-\beta H(\omega)},
    \qquad A \in \mathcal{F}.
  \]
  Assim, estados de menor energia recebem maior peso, enquanto a função de
  partição normaliza esses pesos para que a probabilidade total seja 1.
\end{example}

\begin{exercise}\label{exercise:gibbs-high-temperature}
  No contexto do \Cref{ex:gibbs-measure}, fixe um evento \(A \in \mathcal{F}\).
  Calcule o limite de alta temperatura
  \[
    \lim_{\beta \downarrow 0} \mathbb{P}_\beta(A).
  \]
  Interprete o resultado em termos da medida uniforme sobre \(\Omega\).
\end{exercise}

\begin{exercise}\label{exercise:gibbs-low-temperature}
  No contexto do \Cref{ex:gibbs-measure}, defina
  \[
    h_{\min} = \min_{\omega \in \Omega} H(\omega),
    \qquad
    \Omega_{\min} = \{\omega \in \Omega : H(\omega) = h_{\min}\}.
  \]
  Calcule o limite de baixa temperatura
  \[
    \lim_{\beta \to \infty} \mathbb{P}_\beta(A).
  \]
  Mostre que, no limite, a medida se torna uniforme sobre os estados de menor
  energia. Em particular, descreva o que acontece quando o estado de menor
  energia é único.
\end{exercise}
